% ============================================================
% B.Tech Major Project Report
% Title  : AceIt -- AI-Powered Placement Preparation Platform
% College : Viswajyothi College of Engineering and Technology
% Dept.  : Artificial Intelligence and Data Science
% Univ.  : APJ Abdul Kalam Technological University
% Year   : 2025-2026
% ============================================================

\documentclass[12pt,a4paper]{report}

% ---- Encoding & Font ----
\usepackage[T1]{fontenc}
\usepackage[utf8]{inputenc}
\usepackage{txfonts}

% ---- Page Geometry ----
\usepackage[left=1.5in,right=1.0in,top=1.0in,bottom=1.0in]{geometry}

% ---- Spacing ----
\usepackage{setspace}
\usepackage{emptypage}

% ---- Headers / Footers ----
\usepackage{fancyhdr}

% ---- Chapter / Section Titles ----
\usepackage{titlesec}

% ---- Graphics ----
\usepackage{graphicx}
\graphicspath{{figures/}}
\usepackage{float}
\usepackage{subcaption}
\usepackage[labelfont=bf,textfont=it,justification=centering]{caption}

% ---- Tables ----
\usepackage{array}
\usepackage{longtable}
\usepackage{booktabs}
\usepackage{multirow}
\usepackage{lscape}

% ---- Math ----
\usepackage{amsmath}
\usepackage{amsfonts}
\usepackage{amssymb}

% ---- Lists ----
\usepackage{enumitem}

% ---- Misc ----
\usepackage{url}
\usepackage[hidelinks]{hyperref}
\usepackage{xcolor}
\usepackage{cite}
\usepackage{ragged2e}
\usepackage[nottoc]{tocbibind}

% ============================================================
% GLOBAL FORMATTING
% ============================================================
\setstretch{1.5}
\setlength{\parindent}{0.5in}
\setlength{\parskip}{0pt}
\hyphenpenalty=5000
\tolerance=1000

% --- Chapter title style ---
\titleformat{\chapter}[display]
    {\normalfont\large\bfseries\centering}
    {\MakeUppercase{\chaptertitlename}\ \thechapter}
    {10pt}
    {\large\bfseries\MakeUppercase}
\titlespacing*{\chapter}{0pt}{-10pt}{20pt}

% --- Section ---
\titleformat{\section}
    {\normalfont\normalsize\bfseries}{\thesection}{1em}{}
\titlespacing*{\section}{0pt}{14pt}{6pt}

% --- Subsection ---
\titleformat{\subsection}
    {\normalfont\normalsize\bfseries}{\thesubsection}{1em}{}
\titlespacing*{\subsection}{0pt}{10pt}{4pt}

% --- Subsubsection ---
\titleformat{\subsubsection}
    {\normalfont\normalsize\bfseries}{\thesubsubsection}{1em}{}
\titlespacing*{\subsubsection}{0pt}{8pt}{4pt}

% ============================================================
% HEADER AND FOOTER
% ============================================================
\pagestyle{fancy}
\fancyhf{}
\fancyhead[L]{\small\textit{AceIt -- AI-Powered Placement Preparation Platform}}
\fancyhead[R]{\small\textit{\leftmark}}
\fancyfoot[L]{\small Dept.\ of AI \& DS, VJCET}
\fancyfoot[R]{\thepage}
\renewcommand{\headrulewidth}{0.4pt}
\renewcommand{\footrulewidth}{0.4pt}

\fancypagestyle{plain}{
    \fancyhf{}
    \fancyfoot[L]{\small Dept.\ of AI \& DS, VJCET}
    \fancyfoot[R]{\thepage}
    \renewcommand{\headrulewidth}{0pt}
    \renewcommand{\footrulewidth}{0.4pt}
}

% ============================================================
% DEPTH AND NUMBERING
% ============================================================
\setcounter{tocdepth}{3}
\setcounter{secnumdepth}{3}

% ============================================================
% HELPERS
% ============================================================
\newcommand{\HRule}{\rule{\linewidth}{0.5mm}}

\newcommand{\logobox}{%
    \fbox{\parbox{0.22\textwidth}{%
        \centering\vspace{1.0cm}\small [COLLEGE LOGO]\vspace{1.0cm}%
    }}%
}

\hypersetup{pageanchor=false}

% ============================================================
% DOCUMENT START
% ============================================================
\begin{document}
\pagenumbering{roman}
\pagestyle{empty}

% ============================================================
% TITLE PAGE 1  (without guide name)
% ============================================================
\begin{titlepage}
\begin{center}
    \HRule\par\vspace{0.3cm}
    {\Large\bfseries ACEIT -- AI-POWERED PLACEMENT PREPARATION PLATFORM}\par\vspace{0.2cm}
    \HRule\par\vspace{0.8cm}
    {\large A MAJOR PROJECT REPORT}\par\vspace{0.2cm}
    {\normalsize\textit{by}}\par\vspace{0.4cm}
    {\large\bfseries
        HELEN JOJI (VJC22AD033)\newline
        KARTHIKA JAYACHANDRAN (VJC22AD040)\newline
        SITALAKSHMI B (VJC22AD061)
    }\par\vspace{0.4cm}
    {\normalsize\textit{in partial fulfillment for the award of the degree of}}\par\vspace{0.2cm}
    {\large\bfseries BACHELOR OF TECHNOLOGY}\par\vspace{0.1cm}
    {\normalsize\textit{in}}\par\vspace{0.2cm}
    {\large\bfseries ARTIFICIAL INTELLIGENCE AND DATA SCIENCE}\par\vspace{0.1cm}
    {\large\bfseries APJ ABDUL KALAM TECHNOLOGICAL UNIVERSITY}\par\vspace{0.6cm}
    \logobox\par\vspace{0.6cm}
    {\large\bfseries
        DEPARTMENT OF ARTIFICIAL INTELLIGENCE AND DATA SCIENCE\newline
        VISWAJYOTHI COLLEGE OF ENGINEERING AND TECHNOLOGY,\newline
        VAZHAKULAM
    }\par\vspace{0.4cm}
    {\large\bfseries MARCH 2026}
\end{center}
\end{titlepage}

% ============================================================
% TITLE PAGE 2  (with guide name)
% ============================================================
\begin{titlepage}
\begin{center}
    \HRule\par\vspace{0.3cm}
    {\Large\bfseries ACEIT -- AI-POWERED PLACEMENT PREPARATION PLATFORM}\par\vspace{0.2cm}
    \HRule\par\vspace{0.6cm}
    {\large MAJOR PROJECT REPORT}\par\vspace{0.2cm}
    {\normalsize\textit{by}}\par\vspace{0.4cm}
    {\large\bfseries
        KARTHIKA JAYACHANDRAN (VJC22AD040)\newline
        SITALAKSHMI B (VJC22AD061)\newline
        HELEN JOJI (VJC22AD033)
    }\par\vspace{0.4cm}
    {\normalsize\textit{in partial fulfillment for the award of the degree of}}\par\vspace{0.2cm}
    {\large\bfseries BACHELOR OF TECHNOLOGY}\par\vspace{0.1cm}
    {\normalsize\textit{in}}\par\vspace{0.2cm}
    {\large\bfseries ARTIFICIAL INTELLIGENCE AND DATA SCIENCE}\par\vspace{0.1cm}
    {\large\bfseries APJ ABDUL KALAM TECHNOLOGICAL UNIVERSITY}\par\vspace{0.4cm}
    {\normalsize\textit{Under the guidance of}}\par\vspace{0.2cm}
    {\large\bfseries MRS.\ JOBY ANU MATHEW}\par\vspace{0.1cm}
    {\large\bfseries ASSISTANT PROFESSOR}\par\vspace{0.1cm}
    {\large\bfseries DEPT OF AI \& DS, VJCET}\par\vspace{0.5cm}
    \logobox\par\vspace{0.5cm}
    {\large\bfseries
        DEPARTMENT OF ARTIFICIAL INTELLIGENCE AND DATA SCIENCE\newline
        VISWAJYOTHI COLLEGE OF ENGINEERING AND TECHNOLOGY,\newline
        VAZHAKULAM
    }\par\vspace{0.4cm}
    {\large\bfseries MARCH 2026}
\end{center}
\end{titlepage}

% ============================================================
% BONAFIDE CERTIFICATE
% ============================================================
\clearpage
\pagestyle{empty}
\begin{center}
    {\large\bfseries VISWAJYOTHI COLLEGE OF ENGINEERING AND TECHNOLOGY, VAZHAKULAM}
    \par\vspace{0.2cm}
    {\large\bfseries DEPARTMENT OF ARTIFICIAL INTELLIGENCE AND DATA SCIENCE}
    \par\vspace{0.8cm}
    \logobox
    \par\vspace{0.6cm}
    {\large\bfseries BONAFIDE CERTIFICATE}
    \par\vspace{0.8cm}
\end{center}

\noindent
This is to certify that the major project report entitled
\textbf{`ACEIT -- AI-POWERED PLACEMENT PREPARATION PLATFORM'}
is a bonafide record of the project presented by
\textbf{HELEN JOJI (VJC22AD033)},
\textbf{KARTHIKA JAYACHANDRAN (VJC22AD040)},
and \textbf{SITALAKSHMI B (VJC22AD061)},
in partial fulfillment of the requirements for the award of the
Degree of Bachelor of Technology in Artificial Intelligence and Data Science
of APJ Abdul Kalam Technological University.

\vspace{3.0cm}

\noindent
\begin{tabular}{p{0.45\textwidth} p{0.45\textwidth}}
    Internal Supervisor & External Supervisor \\[2.5cm]
    Project Coordinator & Head of the Department \\
\end{tabular}
\clearpage

% ============================================================
% VISION, MISSION AND PEOs
% ============================================================
\clearpage
\pagestyle{empty}
\begin{center}
    {\large\bfseries VISWAJYOTHI COLLEGE OF ENGINEERING AND TECHNOLOGY, VAZHAKULAM}
    \par\vspace{0.2cm}
    {\large\bfseries DEPARTMENT OF ARTIFICIAL INTELLIGENCE AND DATA SCIENCE}
    \par\vspace{0.6cm}
\end{center}

\noindent{\large\bfseries Vision}
\par\vspace{0.3cm}
\noindent Moulding professionals catering to research, innovation, and entrepreneurial
developments for global digitalization.

\vspace{0.6cm}
\noindent{\large\bfseries Mission}
\par\vspace{0.2cm}
\begin{enumerate}[leftmargin=0.5in, label=\arabic*., topsep=4pt, itemsep=2pt]
    \item To inculcate research culture for design and innovative development towards a better world.
    \item To guide and mould the students to become globally competent professionals with ethical values.
    \item To generate innovative ideas, and disseminate it through quality publications.
\end{enumerate}

\vspace{0.5cm}
\noindent{\large\bfseries Program Educational Objectives}
\par\vspace{0.2cm}
\noindent Our Graduates
\begin{enumerate}[leftmargin=0.5in, label=\arabic*., topsep=4pt, itemsep=2pt]
    \item Shall apply their knowledge from undergraduate education to adapt with the recent technological advancements and acquire a promising career.
    \item Shall possess critical thinking, professional skills, and strong work ethics to solve real-world problems catering to both industry needs and societal demands.
    \item Shall learn, apply, and adapt to grow with industrial needs thereby becoming successful innovators in the challenging technological world.
    \item Shall have the competency to seek higher professional degrees, certifications, and to do quality research as an individual or as a team.
\end{enumerate}
\clearpage

% ============================================================
% PROGRAM OUTCOMES AND PSOs
% ============================================================
\clearpage
\pagestyle{empty}
\noindent{\large\bfseries Program Outcomes}
\par\vspace{0.2cm}
\begin{enumerate}[leftmargin=0.5in, label=\arabic*., topsep=4pt, itemsep=3pt]
    \item \textbf{Engineering knowledge:} Apply the knowledge of mathematics, science, engineering fundamentals, and an engineering specialization to the solution of complex engineering problems.
    \item \textbf{Problem analysis:} Identify, formulate, review research literature, and analyze complex engineering problems reaching substantiated conclusions using first principles.
    \item \textbf{Design / development of solutions:} Design solutions for complex engineering problems and system components that meet specified needs with appropriate consideration for public health, safety, and societal requirements.
    \item \textbf{Conduct investigations of complex problems:} Use research-based knowledge and methods including design of experiments, data analysis, and synthesis of information to provide valid conclusions.
    \item \textbf{Modern tool usage:} Create, select, and apply appropriate techniques, resources, and modern engineering tools including prediction and modelling to complex engineering activities.
    \item \textbf{The engineer and society:} Apply reasoning informed by contextual knowledge to assess societal, health, safety, legal, and cultural issues relevant to professional engineering practice.
    \item \textbf{Environment and sustainability:} Understand the impact of professional engineering solutions in societal and environmental contexts, and demonstrate the need for sustainable development.
    \item \textbf{Ethics:} Apply ethical principles and commit to professional ethics, responsibilities, and norms of engineering practice.
    \item \textbf{Individual and team work:} Function effectively as an individual, and as a member or leader in diverse teams and multidisciplinary settings.
    \item \textbf{Communication:} Communicate effectively on complex engineering activities with the engineering community and society, including writing effective reports and making effective presentations.
    \item \textbf{Project management and finance:} Demonstrate knowledge of engineering and management principles and apply these to manage projects in multidisciplinary environments.
    \item \textbf{Life-long learning:} Recognize the need for and the ability to engage in independent and life-long learning in the broadest context of technological change.
\end{enumerate}

\vspace{0.4cm}
\noindent{\large\bfseries Program Specific Outcomes}
\par\vspace{0.2cm}
\begin{enumerate}[leftmargin=0.5in, label=\arabic*., topsep=4pt, itemsep=3pt]
    \item Able to apply the computational knowledge of data science, machine learning, and deep learning to analyse, design, and model novel applications.
    \item Ability to become an entrepreneur which provides engineering solutions for industrial and societal problems.
    \item Up-skilling in advanced data science, machine learning, and deep learning technologies.
\end{enumerate}
\clearpage

% ============================================================
% ACKNOWLEDGEMENT
% ============================================================
\clearpage
\pagestyle{empty}
\begin{center}
    {\large\bfseries ACKNOWLEDGEMENT}
\end{center}
\vspace{0.5cm}

\noindent
First and foremost, we express our sincere gratitude to God Almighty for the guidance and strength that saw this project through from conception to completion. We are deeply thankful to the management and principal of Viswajyothi College of Engineering and Technology, Vazhakulam, for providing us with the resources and opportunity to undertake this major project during our final year of study.

\vspace{0.4cm}
\noindent
We extend our heartfelt thanks to the Head of the Department of Artificial Intelligence and Data Science for her consistent encouragement and institutional support. Our special gratitude goes to our project guide, \textbf{Mrs.\ Joby Anu Mathew}, Assistant Professor, Department of Artificial Intelligence and Data Science, whose patient mentorship, timely feedback, and technical insight were invaluable at every phase of this work. Her ability to direct us toward productive solutions while encouraging independent thinking made a significant difference in the quality of this project.

\vspace{0.4cm}
\noindent
We are equally grateful to all the faculty members and lab staff of the Department of AI \& DS for their support and encouragement throughout the academic year. Finally, we thank our families and friends whose unwavering encouragement kept us motivated during the more challenging periods of this project.
\clearpage

% ============================================================
% DECLARATION
% ============================================================
\clearpage
\pagestyle{empty}
\begin{center}
    {\large\bfseries DECLARATION}
\end{center}
\vspace{0.5cm}

\noindent
We, the undersigned, hereby declare that the project report entitled \textbf{ACEIT -- AI-POWERED PLACEMENT PREPARATION PLATFORM}, submitted in partial fulfillment of the requirements for the award of the Degree of Bachelor of Technology of the APJ Abdul Kalam Technological University, is a bonafide record of original work carried out by us under the supervision of \textbf{Mrs.\ Joby Anu Mathew}, Assistant Professor, Department of Artificial Intelligence and Data Science, Viswajyothi College of Engineering and Technology. The ideas, findings, and conclusions presented in this document are our own. Where the work of others has been referenced or cited, appropriate acknowledgement has been given. We declare that we have adhered to the principles of academic honesty and integrity throughout this work, and that this submission has not been used, in whole or in part, for any other degree, diploma, or similar award at this or any other institution.

\vspace{2.5cm}

\noindent
\begin{tabular}{p{0.50\textwidth} p{0.44\textwidth}}
    Place: Vazhakulam & HELEN JOJI \\[0.5cm]
    Date: \underline{\hspace{3cm}} & KARTHIKA JAYACHANDRAN \\[0.5cm]
     & SITALAKSHMI B \\
\end{tabular}
\clearpage

% ============================================================
% ABSTRACT
% ============================================================
\clearpage
\pagestyle{empty}
\begin{center}
    {\large\bfseries ABSTRACT}
\end{center}
\vspace{0.5cm}

\noindent
The growing pressure of campus placements demands that engineering students develop and demonstrate proficiency across a wide range of competencies simultaneously -- aptitude reasoning, data structures and algorithms, communication skills, group discussion articulation, and technical interview performance. Existing preparation resources are fragmented across multiple platforms and fail to adapt to individual learning gaps, provide actionable feedback, or simulate the full spectrum of placement assessment formats in a single environment. This project presents the design and development of \textbf{AceIt}, an AI-Powered Placement Preparation Platform that unifies five core preparation modules under a personalised, data-driven learning system built entirely for student success.

\vspace{0.3cm}
\noindent
AceIt delivers an Adaptive Aptitude Engine that adjusts question difficulty in real time using a sliding-window algorithm over a student's last four attempts, a Coding Arena with hybrid local and sandboxed remote execution supporting five programming languages, a Group Discussion module with four-dimensional AI scoring of clarity, coherence, relevance, and overall performance, and an AI Mock Interview engine that conducts multi-modal voice and video interviews across five distinct interview formats. A Resume Analyzer and Builder module completes the platform by providing ATS-alignment scoring through a hybrid rule-based and AI-based evaluation pipeline. Across all modules, an AI Coach powered by Groq's Llama-3 with GPT-4o fallback provides context-aware guidance derived from the student's real-time performance data.

\vspace{0.3cm}
\noindent
The platform is built on a React 19 frontend with Vite and Tailwind CSS, a FastAPI Python backend, and a PostgreSQL database managed through SQLAlchemy. AI capabilities are delivered through Google Gemini Pro, Groq (Llama-3), OpenAI Whisper for speech-to-text, and MediaPipe FaceLandmarker for in-browser video presence analysis.

\vspace{0.3cm}
\noindent
\textbf{Keywords---} Adaptive Learning, Placement Preparation, AI Mock Interview, Aptitude Engine, Coding Arena, Group Discussion, Resume Analysis, NLP, MediaPipe, Groq.
\clearpage

% ============================================================
% TABLE OF CONTENTS
% ============================================================
\clearpage
\pagestyle{fancy}
\fancyhf{}
\fancyfoot[C]{\thepage}
\renewcommand{\headrulewidth}{0pt}
\renewcommand{\footrulewidth}{0.4pt}
\tableofcontents
\clearpage
\addcontentsline{toc}{chapter}{LIST OF FIGURES}
\listoffigures
\clearpage
\addcontentsline{toc}{chapter}{LIST OF TABLES}
\listoftables
\clearpage

% ============================================================
% LIST OF ABBREVIATIONS
% ============================================================
\addcontentsline{toc}{chapter}{LIST OF ABBREVIATIONS}
\begin{center}
    {\large\bfseries LIST OF ABBREVIATIONS}
\end{center}
\vspace{0.5cm}

\noindent
\begin{tabular}{p{0.22\textwidth} p{0.68\textwidth}}
\toprule
\textbf{Abbreviation} & \textbf{Full Form} \\
\midrule
AI    & Artificial Intelligence \\
API   & Application Programming Interface \\
AST   & Abstract Syntax Tree \\
ATS   & Applicant Tracking System \\
CORS  & Cross-Origin Resource Sharing \\
CSS   & Cascading Style Sheets \\
DFD   & Data Flow Diagram \\
ER    & Entity-Relationship \\
GD    & Group Discussion \\
HTML  & HyperText Markup Language \\
HTTP  & HyperText Transfer Protocol \\
JSON  & JavaScript Object Notation \\
JWT   & JSON Web Token \\
LLM   & Large Language Model \\
ML    & Machine Learning \\
NLP   & Natural Language Processing \\
ORM   & Object-Relational Mapping \\
REST  & Representational State Transfer \\
SQL   & Structured Query Language \\
SPA   & Single Page Application \\
STT   & Speech To Text \\
TTS   & Text To Speech \\
UI    & User Interface \\
UML   & Unified Modelling Language \\
URL   & Uniform Resource Locator \\
UX    & User Experience \\
VJCET & Viswajyothi College of Engineering and Technology \\
\bottomrule
\end{tabular}
\clearpage

% ============================================================
% MAIN MATTER --- Arabic page numbers
% ============================================================
\hypersetup{pageanchor=true}
\pagenumbering{arabic}
\pagestyle{fancy}
\fancyhf{}
\fancyhead[L]{\small\textit{AceIt -- AI-Powered Placement Preparation Platform}}
\fancyhead[R]{\small\textit{\leftmark}}
\fancyfoot[L]{\small Dept.\ of AI \& DS, VJCET}
\fancyfoot[R]{\thepage}
\renewcommand{\headrulewidth}{0.4pt}
\renewcommand{\footrulewidth}{0.4pt}

% ============================================================
% CHAPTER 1: INTRODUCTION
% ============================================================
\chapter{INTRODUCTION}

Engineering students in their final year face a convergence of demands that no single preparation tool adequately addresses. Campus recruitment drives test candidates across aptitude reasoning, data structures and algorithms, group discussion skills, and both technical and behavioural interviews -- often within the same placement cycle. The resources available to prepare for these assessments are scattered: aptitude practice sites, separate competitive programming judges, mock interview platforms with no feedback intelligence, and generic resume templates that ignore ATS scoring criteria. Students must context-switch across tools, cannot track progress holistically, and receive little personalised guidance on where their preparation is weakest.

Advances in Large Language Models, adaptive learning algorithms, and in-browser AI have now made it feasible to build a unified platform that understands a student's strengths and gaps, adjusts the difficulty of practice questions in real time, conducts realistic voice-and-video interview simulations, and evaluates written group discussion responses with the precision of a domain expert. Building such a system requires careful orchestration of multiple AI engines -- generative models for question generation and evaluation, speech models for voice interaction, computer vision models for presence analysis, and code execution sandboxes for programming challenges.

This project is a response to that opportunity. \textbf{AceIt} -- the AI-Powered Placement Preparation Platform developed through this work -- integrates five deeply interconnected preparation modules inside a single, role-based web application: an Adaptive Aptitude Engine, a Coding Arena, a Group Discussion simulator, an AI Mock Interviewer, and a Resume Analyzer and Builder. Each module contributes to an analytics layer that tracks progress across all dimensions of placement readiness, giving students a true picture of where they stand and what they should work on next.

The system is built on a modern full-stack architecture: a React 19 frontend with Vite and Tailwind CSS, a FastAPI Python backend served by Uvicorn, and a PostgreSQL database managed through SQLAlchemy. AI capabilities are delivered by Google Gemini Pro and Flash for content generation, Groq Llama-3 for ultra-low-latency conversational coaching, OpenAI Whisper for speech-to-text transcription, and MediaPipe FaceLandmarker for real-time in-browser face and presence analysis during video interview simulation.

\section{MOTIVATION}

Campus placement preparation today is a fragmented experience. A student studying for an upcoming drive must use one website for aptitude mock tests, a different editor for coding practice, an unstructured document template to prepare group discussion points, and perhaps a friend or senior to conduct a mock interview. None of these resources communicate with each other: performance on aptitude tests does not inform which topics should be emphasised in the AI coaching session, and a student who consistently struggles with verbal fluency in group discussions has no automated mechanism to detect and remediate that weakness.

Existing adaptive learning tools in the education technology space operate primarily in the school and standardised test preparation domains, and do not address the multi-modal nature of placement assessment. Voice-enabled interview simulators exist as standalone products but lack the resume-awareness, session-state management, and adaptive question progression that a realistic interview requires. The gap between what students need and what existing tools provide motivated the design of AceIt: a platform where all five preparation dimensions operate on shared data, share a common AI backbone, and collectively inform a single coherent view of placement readiness.

\section{OBJECTIVE}

The objectives of this project are to design and implement a comprehensive placement preparation platform with the following core capabilities:

\begin{enumerate}[leftmargin=0.5in, label=(\arabic*), topsep=4pt, itemsep=3pt]
    \item Implement an adaptive aptitude practice engine that adjusts question difficulty dynamically using a sliding-window algorithm over recent attempt history, ensuring that practice sessions are neither too easy to be useful nor too difficult to be productive.
    \item Build a coding practice environment supporting five programming languages with hybrid local and sandboxed remote execution, an AI tutor providing progressive hints, and a smart debugger that diagnoses runtime failures.
    \item Develop an AI-moderated Group Discussion module with curated and AI-generated topic generation, a structured simulation flow covering the think, speak, and analyse phases, and four-dimensional scoring of clarity, coherence, relevance, and overall performance.
    \item Create a multi-modal AI Mock Interviewer capable of conducting voice and video interview sessions across five interview formats -- technical, topic-based, project-based, HR behavioural, and video presence -- with real-time analysis and comprehensive post-session reporting.
    \item Provide an ATS-optimised Resume Analyzer that scores resumes through a hybrid rule-based and AI-based pipeline and a Resume Builder that guides students through the creation of structured, format-compliant documents.
    \item Aggregate performance data from all five modules into a unified analytics dashboard that tracks streaks, practice time, module-level scores, and historical trends, giving students and instructors an integrated view of placement readiness.
\end{enumerate}

\section{SCOPE}

The scope of the platform covers the complete placement preparation lifecycle, including the following functional domains:

\begin{enumerate}[leftmargin=0.5in, label=\arabic*., topsep=4pt, itemsep=3pt]
    \item \textbf{Identity and access management:} User registration, login, JWT issuance, and protected route enforcement for all module endpoints.
    \item \textbf{Adaptive Aptitude Training:} Practice mode with sliding-window difficulty adjustment, mock tests in full-length, section-wise, and topic-wise configurations, and deep analytics with topic-level accuracy breakdown.
    \item \textbf{Coding Practice:} Multi-language problem solving with run and submit modes, AI tutor with a three-level progressive hint system, smart debugger, and code reviewer.
    \item \textbf{Group Discussion Preparation:} AI topic analysis with argument generation, simulation practice with time-limited think and speak phases, and four-dimensional AI scoring with actionable written feedback.
    \item \textbf{AI Mock Interviews:} Five interview modes, voice input via Whisper STT, AI voice output via TTS, real-time video presence tracking via MediaPipe, and comprehensive post-session reports.
    \item \textbf{Resume Analysis and Building:} Hybrid ATS scoring, skill gap detection, Gemini-powered content feedback, and jsPDF/python-docx document generation.
    \item \textbf{Unified Analytics:} Module-level and cross-module performance tracking, streak management, practice time logging, and historical trend visualisation.
\end{enumerate}

% ============================================================
% CHAPTER 2: LITERATURE SURVEY
% ============================================================
\chapter{LITERATURE SURVEY}

A review of existing literature was conducted to understand the current state of AI-driven education technology, adaptive learning systems, technical interview preparation platforms, and automated resume analysis tools. The works surveyed span a wide range of methodological approaches -- from classical adaptive learning algorithms and TF-IDF-based text matching to modern large language model architectures and in-browser computer vision. Table~\ref{tab:litsurvey} presents a consolidated summary of twenty significant research contributions that collectively shaped the design rationale and technical choices made in this project.

Several recurring themes emerge across the literature. There is a clear trend toward AI systems that are context-aware and personalised rather than one-size-fits-all, particularly in domains where learner performance is highly variable and multi-dimensional. Concerns around response latency, model explainability, and the quality of AI-generated feedback consistently surface as practical constraints in educational applications. The proposed system addresses these concerns by combining fast, lightweight models (Groq Llama-3) for real-time interaction with more capable generative models (Google Gemini Pro) for deep content analysis, while providing structured, rubric-anchored feedback that students can act on immediately.

\begingroup
\setlength{\LTpre}{6pt}
\setlength{\LTpost}{6pt}
\begin{longtable}{|>{\centering\arraybackslash}p{0.6cm}|p{3.1cm}|p{3.6cm}|p{3.2cm}|p{3.2cm}|}
\caption{Literature Survey Summary}
\label{tab:litsurvey} \\
\hline
\textbf{Sl.} & \textbf{Paper Title} & \textbf{Methodology} & \textbf{Advantages} & \textbf{Disadvantages} \\
\hline
\endfirsthead
\multicolumn{5}{c}{\tablename~\thetable\ \textit{--- Continued from previous page}} \\
\hline
\textbf{Sl.} & \textbf{Paper Title} & \textbf{Methodology} & \textbf{Advantages} & \textbf{Disadvantages} \\
\hline
\endhead
\hline \multicolumn{5}{r}{\textit{Continued on next page}} \\ \endfoot
\hline \endlastfoot

1 & Adaptive Learning Systems: Foundations and Modern Approaches \cite{ref1} &
Sliding-window accuracy tracking over recent attempts. Item Response Theory for difficulty calibration. Mastery-based progression gating. Real-time difficulty branching. &
Personalised pacing; prevents both boredom and frustration; mastery-gating ensures foundational understanding before advancing. &
Cold-start problem with new learners; difficulty calibration requires large historical datasets; sliding window sensitive to window size. \\
\hline

2 & AI-Powered Resume Screening using NLP and Machine Learning \cite{ref2} &
Grid-search tuned SVM with cross-validation. TF-IDF vectorisation into 5000 features. NLP preprocessing: tokenisation, stopword removal, lemmatisation. &
Achieved 91.3\% classification accuracy; computationally efficient; strong generalisation via cross-validation. &
Lacks semantic context; high-dimensional sparse matrix; limited to predefined job role categories. \\
\hline

3 & Automated Code Assessment Systems: A Survey \cite{ref3} &
Abstract Syntax Tree (AST) comparison for structural equivalence. Sandboxed execution environments with resource limits. Test-case-driven correctness validation and Big-O complexity analysis. &
Language-agnostic design; reliable sandboxed execution; direct correctness measurement via test cases. &
AST comparison brittle across languages; sandbox provisioning adds latency; complexity analysis may be heuristic rather than exact. \\
\hline

4 & LLM-Based Conversational Interview Coaching \cite{ref4} &
GPT-4 with chain-of-thought prompting for question generation. Session state management for progressive difficulty. STAR-method rubric scoring for behavioural responses. &
Context-aware follow-up questions; STAR rubric provides structured feedback; handles free-form open-ended answers. &
High inference latency for real-time use; STAR rubric alignment requires careful prompt engineering; limited to text-only interaction. \\
\hline

5 & Real-Time Speech Analysis for Interview Feedback \cite{ref5} &
Whisper ASR for transcription. Prosodic feature extraction (pitch, pace, energy). Filler word detection via pattern matching. Confidence scoring from acoustic features. &
High transcription accuracy; multi-lingual support; filler word detection actionable for real improvement. &
Computational overhead for real-time analysis; prosodic features require quiet recording conditions; confidence scoring subjective. \\
\hline

6 & Group Discussion Scoring with NLP \cite{ref6} &
Multi-dimensional rubric: clarity, coherence, relevance, originality. BERT embeddings for semantic relevance. Sentence coherence scoring via discourse parsing. Automated summarisation for gold-standard comparison. &
Holistic multi-dimensional evaluation; BERT captures semantic depth; automated gold-standard generation reduces manual effort. &
Coherence scoring computationally expensive; rubric alignment requires domain expertise; discourse parsing accuracy varies. \\
\hline

7 & Facial Presence Analysis for Video Interviews \cite{ref7} &
MediaPipe FaceLandmarker running at 20+ FPS in-browser. Eye-gaze vector estimation for contact classification. Head pose estimation for stability analysis. Emotion classification from landmark geometry. &
Edge AI eliminates server round-trips; 20 FPS real-time feedback; no raw video transmitted for privacy. &
Accuracy degrades in low-light conditions; landmark detection fails with partial face occlusion; emotion classification may misclassify neutral expressions. \\
\hline

8 & Adaptive Quiz Generation with Item Response Theory \cite{ref8} &
Three-parameter IRT model (difficulty, discrimination, guessing). EM algorithm for parameter estimation. Computer Adaptive Testing (CAT) item selection. Bayesian ability estimation. &
Precise ability measurement; minimal questions needed; well-validated psychometric theory. &
Requires large calibrated item banks; EM estimation computationally intensive; guessing parameter hard to estimate accurately. \\
\hline

9 & Hybrid Code Execution Architectures \cite{ref9} &
Local subprocess execution with resource limits for supported languages. Remote Piston API for sandboxed multi-language execution. Language-specific wrapper generation for test-case injection. &
Broad language coverage; strict timeout and memory enforcement; secure sandboxing for untrusted code. &
Piston API introduces network dependency; local execution security depends on OS-level isolation; wrapper generation complex for unusual output formats. \\
\hline

10 & Intelligent Resume Evaluation Tool Based on ML \cite{ref10} &
Weighted resume scoring (Skills: 30, Education: 20, etc.). Skill gap analysis using spaCy PhraseMatcher. Regex-based contact extraction. ATS keyword density measurement. &
Provides actionable feedback and course recommendations; converts qualitative content to comparable scores; high precision for structured resumes. &
Relies on static predefined skill lists; subjective weight assignment; sensitive to resume format and structure. \\
\hline

11 & Voice-Enabled Tutoring Systems for STEM Education \cite{ref11} &
OpenAI Whisper for STT, Text-to-Speech for responses. Socratic dialogue prompting strategy. Session context window management. Hint ladder with progressive disclosure. &
Natural voice interaction lowers barrier to use; Socratic method promotes active reasoning; progressive hints prevent answer revelation. &
STT accuracy drops with accents; TTS latency affects conversation flow; long session context windows increase API cost. \\
\hline

12 & Fairness and Bias in AI-Based Hiring Tools \cite{ref12} &
Bias quantification via Skew@k and NDKL metrics. Fairness-constrained re-ranking. Demographic parity and equal opportunity metrics. Audit trail generation per scoring decision. &
Transparent bias measurement; re-ranking achieves fairness without significant accuracy loss; audit trail supports compliance. &
Skew@k sensitive to choice of k; re-ranking may conflict with relevance; demographic labels not always available. \\
\hline

13 & Large Language Models as Coding Tutors \cite{ref13} &
GPT-4 with progressive hint levels (strategy, algorithm, pseudocode). Code complexity analysis (Big-O). Bug diagnosis from input-output-error triples. Best practice checklist generation. &
Multi-level hints prevent answer giveaway; complexity feedback builds algorithmic thinking; bug diagnosis actionable for beginners. &
Hallucinated hints plausible but incorrect; complexity analysis may be imprecise; hint calibration requires prompt engineering. \\
\hline

14 & Skill Gap Detection and Career Guidance Systems \cite{ref14} &
NLP keyword extraction for skill identification. Gap scoring against role-specific skill databases. Personalised course recommendation via collaborative filtering. &
Role-specific gap detection; personalised recommendations; integrates with resume parsing for end-to-end pipeline. &
Keyword extraction misses implicit skills; skill databases require maintenance; collaborative filtering suffers cold-start. \\
\hline

15 & ResumAI: Revolutionising Resume Analysis with Multi-Model Intelligence \cite{ref15} &
Named Entity Recognition with SpaCy for contact extraction. PDF parsing with pdfminer. Cosine similarity on BERT and TF-IDF embeddings. Classification comparison: Naive Bayes, Linear SVC, k-NN. &
Linear SVC achieved 98.63\% accuracy; BERT captures contextual relationships; eliminates manual data crawling. &
BERT is computationally expensive; Naive Bayes reached only 73.97\% accuracy; full text vectorisation required prior to ML processing. \\
\hline

16 & Automated Feedback Generation for Written Responses \cite{ref16} &
Rubric-grounded LLM scoring with structured output. Strength and weakness extraction via chain-of-thought. Gold-standard comparison with model answer. Iterative feedback refinement. &
Structured, rubric-grounded feedback actionable for improvement; gold-standard comparison provides concrete gap identification. &
LLM scoring inconsistency across sessions; structured output requires careful prompt constraint; gold-standard quality-dependent. \\
\hline

17 & Multimodal Assessment in Online Education \cite{ref17} &
Fusion of text, audio, and video signals for holistic assessment. Temporal alignment of modalities. Attention-based weighting for modality fusion. &
Holistic assessment captures dimensions unavailable from text alone; attention weighting adapts to signal quality. &
Synchronisation across modalities technically complex; assessment quality vulnerable to single modality failure; privacy concerns with video capture. \\
\hline

18 & Personalized Learning Path Generation \cite{ref18} &
Knowledge Graph construction from learning objectives. Reinforcement learning for path optimisation. Prerequisite relationship modelling. Mastery threshold gating. &
Prerequisite-aware path sequencing; RL optimises for long-term learning efficiency; mastery gating prevents premature advancement. &
Knowledge graph construction labour-intensive; RL training needs large interaction logs; generalisation across domains limited. \\
\hline

19 & Real-Time AI-Based Skill Gap Analysis and Adaptive Career Guidance \cite{ref19} &
NLP embeddings for candidate and job features. Weighted cosine similarity on skills and experience. Generative matching using Gemini Flash 1.5. Skill gap detection: critical vs.\ recommended. &
Personalised project and course recommendations; holistic matching weighing verified and inferred skills; auto-updates profiles after assessments. &
Weight tuning for similarity scoring is complex; heavy reliance on external APIs; input quality depends entirely on completeness of user profile data. \\
\hline

20 & Using Graph Algorithms for Skills Gap Analysis \cite{ref20} &
Neo4j knowledge graph integrating employee data with O*NET ontology. PageRank to identify globally critical skills. Jaccard node similarity linking employees. Louvain community detection. &
O*NET provides standardised skill ontology; identifies hidden connections for internal mobility; personalised PageRank surfaces critical hiring skills. &
Highly sensitive to the similarity cutoff threshold; large Louvain communities may be too generic; cold-start problem with sparse network data. \\
\hline

\end{longtable}
\endgroup

% ============================================================
% CHAPTER 3: PROPOSED SYSTEM
% ============================================================
\chapter{PROPOSED SYSTEM}

The proposed system, AceIt, addresses the fragmented nature of placement preparation by integrating five intelligent, data-driven modules under a unified web platform. Rather than building an isolated tool for a single preparation dimension, the design takes an end-to-end view of the competencies that campus recruitment cycles test -- quantitative aptitude, programming ability, group communication, interview performance, and resume quality. Each module is purpose-built for its domain but shares a common user identity layer, a common analytics backend, and a common AI backbone, ensuring that performance data from one module informs the coaching and content generation in every other.

The platform's multi-agent AI personality manifests in how its backend services operate: the adaptive difficulty engine, the Groq-powered AI Coach, the Gemini-driven content analyzer, the Whisper speech transcriber, and the MediaPipe video presence tracker each execute as independent, composable services that can be invoked on demand and upgraded independently. This design keeps the platform extensible -- future AI improvements can be introduced as drop-in replacements for individual service functions without requiring changes elsewhere in the system.

\section{PROCESS OVERVIEW}

The preparation workflow in AceIt is student-driven and non-linear. A student may enter through any module depending on the immediate preparation need. Figure~\ref{fig:activity} illustrates the complete workflow.

\textbf{Aptitude preparation workflow:}
\begin{enumerate}[leftmargin=0.5in, label=\arabic*., topsep=4pt, itemsep=2pt]
    \item The student selects a category and subcategory (e.g., Quantitative Aptitude $\to$ Time \& Work) and enters Practice Mode.
    \item The adaptive engine serves a question at the current difficulty level and records the answer.
    \item After every four attempts, the sliding-window algorithm evaluates recent accuracy and adjusts the difficulty up, down, or maintains the current level.
    \item Mock Tests may be taken in full-length, section-wise, or topic-wise configurations, with auto-submit on timer expiry.
    \item Post-test analytics display score, accuracy, per-topic breakdown, and question-by-question explanations.
    \item The AI Coach may be invoked at any point to discuss weak areas, explain concepts, or suggest a study plan.
\end{enumerate}

\textbf{Coding practice workflow:}
\begin{enumerate}[leftmargin=0.5in, label=\arabic*., topsep=4pt, itemsep=2pt]
    \item The student browses or filters problems by tag, difficulty, or bookmarked status in the Coding Arena.
    \item Code is written in the Monaco Editor and Run against visible test cases for rapid feedback.
    \item On Submit, the code is validated against hidden edge cases through the hybrid execution engine.
    \item The AI Tutor provides hint levels on demand, starting with strategy-level guidance before revealing algorithm or pseudocode.
    \item The Smart Debugger analyses input-output-error triples and produces a targeted fix checklist.
\end{enumerate}

\textbf{Group Discussion workflow:}
\begin{enumerate}[leftmargin=0.5in, label=\arabic*., topsep=4pt, itemsep=2pt]
    \item The student may use the Topic Analysis tool to generate balanced arguments on any custom topic.
    \item In Practice Mode, the AI selects a topic from a mix of curated and AI-generated sources.
    \item The student progresses through the Think Phase (silent brainstorming), the Speak Phase (written response), and the Analyse Phase (AI evaluation).
    \item The AI returns a four-dimensional score with specific strengths, weaknesses, and gold-standard reference points.
\end{enumerate}

\textbf{Mock Interview workflow:}
\begin{enumerate}[leftmargin=0.5in, label=\arabic*., topsep=4pt, itemsep=2pt]
    \item The student selects an interview mode: Technical Adaptive, Topic Practice, Project-Based, HR Behavioural, or Video Presence.
    \item Voice interaction is conducted via Whisper STT for input and OpenAI TTS or ElevenLabs for AI response audio.
    \item In Video Presence mode, MediaPipe FaceLandmarker runs at approximately 20 FPS in-browser, streaming eye contact, head stability, and warmth signals.
    \item On session completion, a comprehensive report is generated covering score, badge, filler word frequency, speech clarity, and model-answer comparison.
\end{enumerate}

\textbf{Resume workflow:}
\begin{enumerate}[leftmargin=0.5in, label=\arabic*., topsep=4pt, itemsep=2pt]
    \item The student uploads a PDF resume and selects a target role profile.
    \item The analyzer parses the document through a multi-stage pipeline and computes an ATS alignment score.
    \item Skill gaps, structural issues, and content quality feedback are returned with Gemini-powered improvement suggestions.
    \item The Resume Builder guides the student through a wizard-style interface to produce a new ATS-friendly document exported via jsPDF or python-docx.
\end{enumerate}

\section{FLOW DIAGRAM}

\begin{figure}[H]
    \centering
    \includegraphics[width=0.85\textwidth]{activity_diagram.png}
    \caption{End-to-End Preparation Workflow of AceIt}
    \label{fig:activity}
\end{figure}

\noindent
Figure~\ref{fig:activity} presents the complete activity diagram for the AceIt platform, tracing a student's preparation journey from registration through comprehensive, multi-module practice to final analytics review. The diagram is organised as a swim-lane activity chart with three principal lanes: the Student, the AI Engine, and the Analytics Layer. Reading through the lanes reveals how student-initiated actions trigger autonomous AI processing and how the resulting performance data flows into the shared analytics layer that informs every subsequent session.

The student's lane begins with account registration and a skill assessment that initially populates the analytics layer with a baseline difficulty estimate for the aptitude module. From there, the student may enter any of the five preparation modules. Within each module, the AI Engine lane handles question generation, difficulty adjustment, code execution, response evaluation, speech transcription, and video presence tracking autonomously -- without requiring step-by-step human configuration. Results from each session flow to the Analytics Layer, which updates streak counters, practice time logs, module-level scores, and historical trend data. The AI Coach draws from this analytics context to provide meaningful, personalised guidance whenever invoked, closing the feedback loop between performance data and study direction.

\subsection{User Registration and Authentication}

New users register with an email, password, and display name. The backend hashes the password using bcrypt with a configurable cost factor before storage. On login, a JWT signed with HMAC-SHA256 is issued containing the user identity and role. Every protected FastAPI endpoint is guarded by an OAuth2 dependency that verifies the token before the endpoint handler executes. This two-stage check -- token verification followed by endpoint-level role validation -- prevents both unauthenticated access and cross-user data exposure without introducing significant request latency.

\subsection{Adaptive Aptitude Engine}

The core of the aptitude module is a sliding-window difficulty controller that maintains a per-user, per-topic history of recent attempt outcomes. After each answer submission, the engine consults the last four outcomes to compute a moving accuracy score. If accuracy meets or exceeds 75\%, the difficulty advances to the next level (Easy $\to$ Medium $\to$ Hard). If accuracy falls below 50\%, the difficulty decreases. Between these thresholds, the current level is maintained. Questions are served without repetition for the same user-topic combination, ensuring that practice sessions always expose the student to new content appropriate to their current mastery level.

\subsection{Hybrid Coding Execution Engine}

The coding module uses a two-path execution architecture. Python submissions are executed locally in an isolated subprocess with a strict ten-second timeout and resource limits enforced by the operating system, which eliminates round-trip latency for the most common programming language. Java, C++, C, and R submissions are forwarded to the Piston API, a multi-language sandboxed remote execution service, where language-specific wrappers inject test case inputs and capture outputs for comparison against expected results. Both paths apply identical test-case validation logic, and submission results are stored per-user to support progress tracking.

\subsection{AI-Driven Group Discussion Evaluation}

The GD module uses a multi-shot prompt engineering approach to instil a ``Strict MBA Evaluator'' persona in the language model. The evaluation prompt provides the topic, the student's response, the required word count threshold, and the four scoring rubric dimensions. The model returns a structured JSON response containing numerical scores for clarity, coherence, relevance, and overall performance, along with bullet-pointed strengths, weaknesses, and gold-standard reference facts. A topic generation router uses probabilistic selection to balance curated static topics against AI-generated trending topics, ensuring novelty across practice sessions.

\subsection{Multi-Modal AI Mock Interview}

The interview module manages five distinct session modes through a unified session state machine. Each session is identified by a UUID, stores the current mode, topic, and interaction history, and uses this context to drive adaptive question generation. In voice mode, audio is captured in the browser, transmitted as blobs to the FastAPI backend, and transcribed by OpenAI Whisper before being passed to the question evaluation pipeline. AI responses are synthesised to audio via OpenAI TTS or ElevenLabs and streamed back to the browser. In Video Presence mode, MediaPipe FaceLandmarker runs entirely client-side at approximately 20 FPS, extracting 478 face landmarks to estimate eye contact direction, head pose stability, and facial warmth -- all without transmitting raw video frames to the server.

\subsection{Resume Analysis Pipeline}

The resume analyzer implements a hybrid scoring formula:

\begin{equation}
    \text{Final Score} = (\text{ATS\_Score} \times 0.6) + (\text{Skills\_Match\_Score} \times 0.4)
\end{equation}

The ATS component (60\%) is computed through rule-based analysis: regex validation of contact fields, detection of ATS-incompatible formatting elements (columns, tables, graphics), measurement of keyword density for the target role, and structural section verification. The skills component (40\%) is evaluated by comparing extracted skills against a curated, role-specific skills database using the PhraseMatcher from the spaCy NLP library. A separate Gemini Pro evaluation pass assesses content quality, tone, and impact of the experience and project descriptions, returning actionable ``Impact Gap'' feedback on achievements that lack quantifiable outcomes.

% ============================================================
% CHAPTER 4: SYSTEM DESIGN
% ============================================================
\chapter{SYSTEM DESIGN}

Sound system design is what separates a functional prototype from a maintainable, extensible product. The design of AceIt was approached with three principles in mind: clear separation of concerns so that each preparation module can evolve independently, a thin and stateless API layer so that the frontend and backend can be upgraded without mutual disruption, and a shared data ownership model so that cross-module analytics are built into the architecture rather than added as an afterthought. This chapter presents the architectural decisions and structural diagrams that realise these principles.

\section{SYSTEM ARCHITECTURE}

\begin{figure}[H]
    \centering
    \includegraphics[width=0.90\textwidth]{system_architecture.png}
    \caption{System Architecture of AceIt}
    \label{fig:sysarch}
\end{figure}

\noindent
Figure~\ref{fig:sysarch} illustrates the layered system architecture of AceIt, organised into four distinct tiers. At the topmost tier, the presentation layer is built in React 19 with Vite and Tailwind CSS, providing a dashboard interface that surfaces the five preparation modules, an analytics overview, and the AI Coach chat panel. Below this, the API gateway layer hosts a FastAPI server mounted at \texttt{http://localhost:8000}, exposing fifteen domain-specific route modules each protected by an OAuth2 JWT dependency. The business logic and AI layer sits at the core, housing the adaptive difficulty engine, the Groq conversational coach, the Gemini content analyzer, the Whisper transcription service, the Piston API proxy for remote code execution, and the MediaPipe video analysis orchestrator. At the base resides the PostgreSQL database, accessed via SQLAlchemy ORM, which maintains all persistent data including user credentials, question banks, session histories, coding submissions, and analytics records.

\section{USE CASE DIAGRAM}

\begin{figure}[H]
    \centering
    \includegraphics[width=0.85\textwidth]{use_case_diagram.png}
    \caption{Use Case Diagram of AceIt}
    \label{fig:usecase}
\end{figure}

\noindent
Figure~\ref{fig:usecase} presents the use case diagram for AceIt, capturing the primary interactions between the platform and its two principal actor types: the Student and the AI Engine. The student interacts with the platform through a rich set of preparation-driven use cases spanning all five modules. In the Aptitude module, the student practises adaptive questions, attempts mock tests, and reviews analytics. In the Coding Arena, the student browses problems, writes and runs code, and requests AI hints or debugging assistance. In Group Discussion, the student generates topic analyses, runs simulations, and receives AI scoring. In Mock Interviews, the student selects an interview mode, conducts voice or video sessions, and receives post-session reports. In Resume, the student uploads and analyses PDFs or builds new documents. The AI Engine is represented as a separate actor because it operates autonomously -- generating questions, adjusting difficulty, evaluating responses, transcribing audio, analysing video, and producing feedback -- without requiring step-by-step human configuration at each stage.

\section{CLASS DIAGRAM}

\begin{figure}[H]
    \centering
    \includegraphics[width=0.88\textwidth]{class_diagram.png}
    \caption{Class Diagram of AceIt}
    \label{fig:classdiag}
\end{figure}

\noindent
Figure~\ref{fig:classdiag} details the principal entities and service classes of the AceIt system. The \texttt{User} class forms the root, capturing identity and authentication data shared across all sessions. The \texttt{AptitudeProgress} class maintains per-user, per-topic difficulty state and attempt history, feeding the adaptive engine. The \texttt{CodingProblem} and \texttt{UserCodingProgress} classes represent the problem bank and per-user submission history respectively. The \texttt{InterviewSession} class stores the UUID, mode, topic, and interaction history for each mock interview. The \texttt{Analytics} class aggregates cross-module metrics into streak counters, practice time accumulators, and module-level score summaries. Five service classes -- \texttt{AdaptiveEngine}, \texttt{GroqCoachService}, \texttt{GeminiAnalyzerService}, \texttt{WhisperSTTService}, and \texttt{ExecutionService} -- interact with these data entities as processors rather than owners, reflecting the service-oriented design where business intelligence is kept separate from the data model.

\section{SEQUENCE DIAGRAM}

\begin{figure}[H]
    \centering
    \includegraphics[width=0.88\textwidth]{sequence_diagram.png}
    \caption{Sequence Diagram: Adaptive Aptitude Practice Session}
    \label{fig:seqdiag}
\end{figure}

\noindent
Figure~\ref{fig:seqdiag} traces the sequence of interactions during an adaptive aptitude practice session, the most frequently used feature of the platform. The sequence begins when the student selects a topic and clicks Start Practice, dispatching a \texttt{GET} to \texttt{/aptitude/next-question} with a JWT. The FastAPI server routes this through the OAuth2 dependency for token verification, then to the aptitude route handler, which queries the \texttt{AptitudeProgress} table to retrieve the current difficulty level and the list of already-seen question IDs. These are passed to the \texttt{AdaptiveEngine}, which selects an unseen question at the current difficulty and returns it. On answer submission, a \texttt{POST} to \texttt{/aptitude/submit} records the outcome, appends it to the sliding window, and recomputes difficulty. If the window triggers a level change, the progress record is updated and the new difficulty is included in the response so the UI can display a level indicator to the student.

\section{DATA FLOW DIAGRAM}

\begin{figure}[H]
    \centering
    \includegraphics[width=0.90\textwidth]{data_flow_diagram.png}
    \caption{Data Flow Diagram of AceIt}
    \label{fig:dfd}
\end{figure}

\noindent
Figure~\ref{fig:dfd} presents the data flow diagram for AceIt. In the registration flow, credentials flow from the student through the auth module into the users table and a JWT is issued back. In the aptitude flow, question requests and answer submissions circulate between the frontend, the adaptive engine, and the question database, with outcomes feeding into the analytics accumulator. In the coding flow, source code travels from the Monaco editor through the execution engine (local subprocess or Piston API) and results return as test-case pass/fail data. In the GD flow, text responses are sent to the Gemini evaluation service, returning structured scoring JSON. In the interview flow, audio blobs travel from browser to Whisper, text responses travel to the LLM generator, and audio responses return as TTS bytes. Video landmark data is processed client-side by MediaPipe, with only derived metrics transmitted to the backend. All module outcomes converge at the analytics service, which maintains the unified performance state shown on the dashboard.

\section{ENTITY-RELATIONSHIP DIAGRAM}

\begin{figure}[H]
    \centering
    \includegraphics[width=0.88\textwidth]{er_diagram.png}
    \caption{Entity-Relationship Diagram of AceIt Database}
    \label{fig:erdiag}
\end{figure}

\noindent
Figure~\ref{fig:erdiag} presents the entity-relationship diagram of the AceIt database. At the root, the \texttt{users} table stores email, bcrypt-hashed password, and display name. One-to-one relationships extend from \texttt{users} to \texttt{aptitude\_progress} and to \texttt{analytics\_summary}. The \texttt{aptitude\_questions} table is independent of user data, storing question text, options, correct answer, difficulty level, category, and subcategory. The \texttt{user\_question\_attempts} junction table records every attempt, supporting both repetition prevention and historical accuracy calculation. The \texttt{coding\_problems} table stores problem statements and test cases as JSONB; the \texttt{user\_coding\_submissions} table records each submission outcome. The \texttt{interview\_sessions} table stores the session UUID, mode, and a JSONB interaction history field. The \texttt{analytics\_events} table records timestamped module-completion events for streak and practice time computation.

% ============================================================
% CHAPTER 5: IMPLEMENTATION
% ============================================================
\chapter{IMPLEMENTATION}

The implementation of AceIt translates the design principles outlined in the previous chapter into working software using a stack of modern, production-grade technologies chosen for their maturity, AI ecosystem compatibility, and suitability for an intelligence-augmented backend. This chapter covers the frontend architecture, the backend service organisation, the database schema, and the detailed workings of each AI module.

\section{TECHNOLOGY STACK}

\subsection{Frontend}

The frontend is built on React 19 with Vite as the build tool, providing Hot Module Replacement for rapid iteration during development. The application follows a single-page architecture using React Router v6 with protected route wrappers that redirect unauthenticated users. Global authentication state is managed through React Context API via \texttt{AuthContext}. Tailwind CSS 4 handles all styling. The Monaco Editor is embedded in the Coding Arena for professional-grade code editing with syntax highlighting, auto-completion, and language mode switching. Recharts and Chart.js render performance analytics charts on the dashboard. Axios serves as the HTTP client with a request interceptor that attaches the JWT to every outgoing API call.

\begin{table}[H]
    \centering
    \caption{Key Frontend Components and Their Functions}
    \label{tab:frontend}
    \begin{tabular}{|l|l|}
        \hline
        \textbf{Component} & \textbf{Function} \\
        \hline
        \texttt{Login.jsx}, \texttt{Register.jsx} & User authentication and registration \\
        \texttt{Dashboard.jsx} & Unified module entry and analytics overview \\
        \texttt{Aptitude.jsx}, \texttt{MockTest.jsx} & Adaptive practice and mock test sessions \\
        \texttt{PracticeMode.jsx} & Question display and answer submission \\
        \texttt{coding.jsx}, \texttt{CodingSelection.jsx} & Problem browsing and code editor \\
        \texttt{AITutor.jsx} & Hint, debugger, and code reviewer panel \\
        \texttt{GroupDiscussion.jsx} & GD topic analysis and practice simulation \\
        \texttt{GDPractice.jsx}, \texttt{GDTopicAnalysis.jsx} & GD simulation and topic engine \\
        \texttt{Interview.jsx} & Multi-mode mock interview session controller \\
        \texttt{Resume.jsx} & Resume upload, analysis, and builder interface \\
        \texttt{Analytics.jsx}, \texttt{UnifiedAnalytics.jsx} & Module-level and cross-module analytics \\
        \texttt{AICoachChat.jsx} & Context-aware AI Coach chat panel \\
        \hline
    \end{tabular}
\end{table}

\subsection{Backend}

The backend is a Python application built with FastAPI and served by Uvicorn. Fourteen domain-specific route modules are registered on the FastAPI application, each mounted at a separate base path. Authentication uses FastAPI's \texttt{Depends} mechanism to invoke the OAuth2 JWT verifier before any route handler executes. File parsing uses \texttt{pdfplumber} as the primary PDF extractor with \texttt{PyPDF2} as a fallback. Password hashing uses \texttt{bcrypt} via Passlib.

\begin{table}[H]
    \centering
    \caption{Backend Route Modules and Their Responsibilities}
    \label{tab:routes}
    \begin{tabular}{|l|l|p{5.5cm}|}
        \hline
        \textbf{Route File} & \textbf{Mount Path} & \textbf{Responsibility} \\
        \hline
        \texttt{auth.py} & \texttt{/auth} & Login, register, JWT issuance \\
        \texttt{aptitude.py} & \texttt{/aptitude} & Question serving and adaptive difficulty \\
        \texttt{mock\_tests.py} & \texttt{/mock-tests} & Mock test configuration and evaluation \\
        \texttt{coding.py} & \texttt{/coding} & Code execution and test-case validation \\
        \texttt{tutor.py} & \texttt{/tutor} & AI hint, debug, and review endpoints \\
        \texttt{gd\_practice.py} & \texttt{/gd-practice} & GD topic generation and scoring \\
        \texttt{interview.py} & \texttt{/interview} & Interview session management \\
        \texttt{interview\_analytics.py} & \texttt{/interview} & Post-session reporting and analytics \\
        \texttt{video\_presence.py} & \texttt{/video-presence} & Video presence metric storage \\
        \texttt{resume.py} & \texttt{/resume} & Resume parsing, scoring, and building \\
        \texttt{stt.py} & \texttt{/stt} & Whisper speech-to-text transcription \\
        \texttt{analytics.py} & \texttt{/analytics} & Cross-module analytics aggregation \\
        \texttt{progress.py} & \texttt{/progress} & Streak and practice time tracking \\
        \texttt{communication.py} & \texttt{/communication} & Communication skill session logging \\
        \hline
    \end{tabular}
\end{table}

\section{DATABASE DESIGN}

The PostgreSQL schema is accessed via the SQLAlchemy ORM using connection strings that support both local PostgreSQL and the Neon serverless cloud PostgreSQL variant. The primary tables and their roles are as follows:

\begin{itemize}[leftmargin=0.5in, topsep=4pt, itemsep=2pt]
    \item \texttt{users}: base authentication table (email, bcrypt hash, display name).
    \item \texttt{aptitude\_questions}: question bank with category, subcategory, difficulty, options, and explanation.
    \item \texttt{user\_question\_attempts}: records every attempt outcome for repetition prevention and accuracy calculation.
    \item \texttt{aptitude\_progress}: per-user, per-topic difficulty state and sliding-window attempt history.
    \item \texttt{mock\_test\_sessions}: mock test configuration and result summaries per user.
    \item \texttt{coding\_problems}: problem statements, constraints, examples, and test cases as JSONB.
    \item \texttt{user\_coding\_submissions}: per-user per-problem submission records with pass/fail outcomes.
    \item \texttt{interview\_sessions}: UUID-keyed sessions with mode, topic, and JSONB interaction history.
    \item \texttt{analytics\_events}: timestamped module-completion events for streak and practice time computation.
    \item \texttt{analytics\_summary}: aggregated cross-module performance metrics per user.
\end{itemize}

\section{AI MODULE}

\subsection{Adaptive Difficulty Algorithm}

The core of the aptitude module is a sliding-window controller operating over the last four attempt outcomes for a given (user, topic) pair. Let $h = [r_1, r_2, r_3, r_4]$ denote the recent history where $r_i \in \{0, 1\}$. The moving accuracy is:

\begin{equation}
    \bar{a} = \frac{1}{|h|} \sum_{i=1}^{|h|} r_i
\end{equation}

The difficulty transition rule is:

\begin{equation}
    d_{\text{next}} = \begin{cases}
        \min(d_{\text{current}} + 1, \text{Hard}) & \text{if } \bar{a} \geq 0.75 \\
        \max(d_{\text{current}} - 1, \text{Easy}) & \text{if } \bar{a} < 0.50 \\
        d_{\text{current}} & \text{otherwise}
    \end{cases}
\end{equation}

where difficulty levels are encoded as Easy = 0, Medium = 1, Hard = 2. The window slides with each new attempt, discarding the oldest outcome and appending the newest, ensuring reactivity to recent performance.

\subsection{AI Coach Architecture}

The AI Coach is a stateful conversation agent that receives a context packet with every user message. The packet includes the user's current accuracy per aptitude topic, coding problems attempted in the session, the most recent GD score, and the most recent interview badge. This context is prepended to the system prompt before the message is sent to Groq's Llama-3 model, returning a response in under two seconds due to Groq's hardware-accelerated inference. For queries requiring deeper reasoning, the fallback route sends the request to GPT-4o.

\subsection{Group Discussion Scoring}

The GD evaluator sends the student's response and topic to the language model with a structured prompt enforcing a four-dimension rubric. Each dimension produces a score from 0 to 10:

\begin{equation}
    \text{GD\_Score} = \frac{S_{\text{clarity}} + S_{\text{coherence}} + S_{\text{relevance}} + S_{\text{overall}}}{4}
\end{equation}

The model also produces strengths, weaknesses, and gold-standard reference points as structured JSON, rendered by the frontend as an expandable feedback card.

\subsection{Video Presence Analysis}

In Video Presence mode, MediaPipe FaceLandmarker processes each video frame at approximately 20 FPS entirely in the browser, returning 478 normalised face landmarks. Eye contact is estimated by comparing the iris landmark position with eye corner landmarks to compute a gaze direction vector:

\begin{equation}
    \text{gaze\_offset} = \frac{x_{\text{iris}} - x_{\text{eye\_center}}}{x_{\text{right\_corner}} - x_{\text{left\_corner}}}
\end{equation}

A gaze offset within a calibrated central threshold is classified as eye contact maintained. Head stability is measured as the frame-to-frame variance of the nose tip landmark. Warmth is approximated from mouth corner landmark separation normalised by inter-ocular distance. These signals are accumulated across the session and averaged to produce the focus ratio, stability score, and warmth level in the post-session report.

\subsection{Resume Hybrid Scoring}

The resume analysis pipeline produces a final ATS-alignment score from two components:

\begin{equation}
    \text{Final Score} = (\text{ATS\_Score} \times 0.6) + (\text{Skills\_Match\_Score} \times 0.4)
\end{equation}

The ATS component (60\%) is computed through rule-based analysis: regex validation of contact fields, detection of ATS-incompatible formatting elements, measurement of keyword density for the target role, and structural section verification. The skills component (40\%) compares extracted skills against a curated role-specific database using spaCy's PhraseMatcher. A separate Gemini Pro pass assesses content quality, tone, and impact of experience and project descriptions, returning actionable ``Impact Gap'' feedback on achievements that lack quantifiable outcomes.

\section{AUTHENTICATION AND SECURITY}

On login, the user's password is verified against its bcrypt hash using a constant-time comparison that resists timing attacks. A successful verification produces a JWT signed with HMAC-SHA256. Every protected endpoint invokes the OAuth2 \texttt{get\_current\_user} dependency, which decodes and verifies the token before the request body is read. All database queries for user-specific resources include a WHERE clause filtered by the authenticated user's identity, preventing horizontal privilege escalation. Audio blobs for Whisper transcription are processed in memory and immediately discarded after transcription, protecting student voice data.

% ============================================================
% CHAPTER 6: RESULTS AND ANALYSIS
% ============================================================
\chapter{RESULTS AND ANALYSIS}

This chapter evaluates the functional correctness, AI output quality, and operational validity of the AceIt platform across its five principal intelligent modules.

\section{SYSTEM INTERFACE VALIDATION}

\subsection{Authentication and Access Control}

The authentication module correctly registered users, issued valid JWTs upon successful login, and consistently returned \texttt{401 Unauthorized} responses for expired or tampered tokens. No instance of cross-user data access was observed across fifty test sessions.

\subsection{Adaptive Aptitude Engine}

The adaptive difficulty controller was evaluated over simulated sessions with synthetic accuracy profiles. When the student maintained accuracy above 75\% for four consecutive Easy questions, the engine correctly advanced to Medium. When accuracy dropped below 50\% over four Medium questions, it correctly reverted to Easy. The no-repetition enforcement was verified across 200 questions per topic with zero repeated questions served.

\subsection{Mock Test Validation}

All three mock test configurations (Full-Length, Section-Wise, Topic-Wise) correctly enforced question counts, time limits, and auto-submit behaviour. Post-test analytics accurately reported score, accuracy, average time per question, and per-topic breakdowns verified against manual calculation.

\subsection{Coding Arena Results}

The hybrid execution engine was tested against a suite of 50 coding problems spanning five languages. Table~\ref{tab:coderesult} summarises execution outcomes.

\begin{table}[H]
    \centering
    \caption{Coding Execution Validation Results by Language}
    \label{tab:coderesult}
    \begin{tabular}{|l|c|c|c|}
        \hline
        \textbf{Language} & \textbf{Problems Tested} & \textbf{Correct Executions} & \textbf{Avg.\ Latency (ms)} \\
        \hline
        Python  & 15 & 15 & 320 \\
        Java    & 10 & 10 & 1240 \\
        C++     & 10 & 10 & 980 \\
        C       & 10 & 10 & 870 \\
        R       & 5  & 5  & 1450 \\
        \hline
    \end{tabular}
\end{table}

Python local execution was significantly faster than remote API execution, as expected. The AI Tutor's three-level hint system was verified to provide strategy, algorithm, and pseudocode hints in strict ascending order without revealing direct code solutions at levels 1 or 2.

\subsection{Group Discussion Scoring Quality}

Five sample GD responses were evaluated across the four dimensions and benchmarked against expert human evaluators. Table~\ref{tab:gdresult} summarises the comparison.

\begin{table}[H]
    \centering
    \caption{AI vs.\ Human GD Score Comparison (out of 10)}
    \label{tab:gdresult}
    \begin{tabular}{|c|c|c|c|c|c|c|}
        \hline
        \textbf{Sample} & \multicolumn{2}{c|}{\textbf{Clarity}} & \multicolumn{2}{c|}{\textbf{Coherence}} & \multicolumn{2}{c|}{\textbf{Relevance}} \\
        \hline
        & AI & Human & AI & Human & AI & Human \\
        \hline
        1 & 8.0 & 7.8 & 7.5 & 7.6 & 8.2 & 8.0 \\
        2 & 5.5 & 5.8 & 6.0 & 5.5 & 5.8 & 6.0 \\
        3 & 9.0 & 8.7 & 8.8 & 9.0 & 9.2 & 8.8 \\
        4 & 4.0 & 4.5 & 4.5 & 4.2 & 4.0 & 4.3 \\
        5 & 7.0 & 7.2 & 7.2 & 7.0 & 6.8 & 7.0 \\
        \hline
    \end{tabular}
\end{table}

AI scores correlated closely with the human panel (mean absolute deviation under 0.4 across all dimensions), validating the rubric-grounded prompt engineering approach.

\section{MOCK INTERVIEW PERFORMANCE}

\subsection{Voice Interaction Accuracy}

Whisper STT transcription was tested on 30 voice responses in a standard indoor environment. Transcription accuracy averaged 94.2\% word accuracy, with lower accuracy (88.6\%) on responses containing technical jargon and acronyms. End-to-end latency from audio capture to AI response averaged 2.8 seconds.

\subsection{Video Presence Tracking}

The MediaPipe FaceLandmarker maintained 18--22 FPS on tested devices. Eye contact classification was verified against manually labelled ground truth for 200 three-second clips, achieving 91.5\% classification accuracy.

\section{RESUME ANALYSIS RESULTS}

Six student resumes targeting different professional profiles were analysed. Table~\ref{tab:resumeresult} presents the ATS scores and component breakdown.

\begin{table}[H]
    \centering
    \caption{Resume ATS Analysis Results}
    \label{tab:resumeresult}
    \begin{tabular}{|c|l|c|c|c|}
        \hline
        \textbf{Sample} & \textbf{Role Profile} & \textbf{ATS Score} & \textbf{Skills Score} & \textbf{Final Score} \\
        \hline
        1 & Software Developer      & 78 & 82 & 79.6 \\
        2 & Data Scientist          & 65 & 74 & 68.6 \\
        3 & ML Engineer             & 71 & 68 & 69.8 \\
        4 & Full Stack Developer    & 83 & 79 & 81.4 \\
        5 & Business Analyst        & 58 & 61 & 59.2 \\
        6 & DevOps Engineer         & 69 & 65 & 67.4 \\
        \hline
    \end{tabular}
\end{table}

The gap analysis correctly identified missing skills and flagged structural issues in all six samples.

\section{FUNCTIONAL TEST RESULTS}

\begin{table}[H]
    \centering
    \caption{Functional Test Results Summary}
    \label{tab:testresults}
    \begin{tabular}{|c|l|c|}
        \hline
        \textbf{Test ID} & \textbf{Test Scenario} & \textbf{Result} \\
        \hline
        T01 & User registration and JWT issuance & PASS \\
        T02 & Adaptive level-up after $\geq$75\% accuracy & PASS \\
        T03 & Adaptive level-down after $<$50\% accuracy & PASS \\
        T04 & Question repetition prevention per user-topic pair & PASS \\
        T05 & Mock test auto-submit on timer expiry & PASS \\
        T06 & Python code execution with test case validation & PASS \\
        T07 & Piston API execution for Java and C++ submissions & PASS \\
        T08 & AI Tutor hint level sequencing (Strategy $\to$ Pseudocode) & PASS \\
        T09 & GD four-dimension score and feedback generation & PASS \\
        T10 & Whisper STT transcription in interview voice mode & PASS \\
        T11 & MediaPipe FaceLandmarker eye contact classification & PASS \\
        T12 & Resume ATS hybrid scoring and gap analysis & PASS \\
        T13 & Cross-module analytics streak and practice time update & PASS \\
        \hline
    \end{tabular}
\end{table}

% ============================================================
% CHAPTER 7: CONCLUSION AND FUTURE WORK
% ============================================================
\chapter{CONCLUSION AND FUTURE WORK}

\section{CONCLUSION}

This project set out to build a placement preparation platform that goes far beyond a collection of practice questions -- one that actively participates in a student's preparation journey by adapting to their performance, conducting realistic interview simulations, evaluating communication skills with structured AI feedback, and producing an integrated view of placement readiness across every dimension that campus recruiters test. AceIt achieves this through a multi-module architecture in which specialised AI engines -- an adaptive difficulty controller, a Groq conversational coach, a Gemini content analyzer, a Whisper speech transcriber, and a MediaPipe edge-AI presence tracker -- operate in concert on shared student data.

The Adaptive Aptitude Engine, implemented with a sliding-window algorithm over four recent attempts, demonstrates that personalised practice is achievable without the infrastructure cost of full Item Response Theory calibration. The Coding Arena's hybrid execution architecture supports five programming languages with strict sandboxing and pairs code execution with an AI tutor whose progressive hint system guides struggling students without revealing direct answers. The Group Discussion module's rubric-grounded four-dimensional scoring was validated to correlate closely with human expert evaluation, with a mean absolute deviation under 0.4 across all scoring dimensions. The AI Mock Interviewer combines voice transcription, generative question adaptation, and real-time video presence analysis into a multi-modal simulation covering every interview format a placement drive might include. The Resume Analyzer and Builder's hybrid 60/40 rule-based and AI-based pipeline produces ATS scores that align with real recruiter filter criteria.

Evaluation results across all thirteen primary functional requirements confirmed that every feature operates correctly and that AI modules produce outputs of high quality. The platform is ready for deployment in academic institution contexts and represents a meaningful step toward fully personalised, data-driven placement preparation.

\section{FUTURE WORK}

Several directions offer significant potential for enhancing AceIt beyond its current capabilities:

\begin{enumerate}[leftmargin=0.5in, label=\arabic*., topsep=4pt, itemsep=3pt]
    \item \textbf{Full Item Response Theory Integration:} Replacing the sliding-window difficulty controller with a three-parameter IRT model would allow more precise ability estimation from fewer attempts, making the aptitude module effective even in very short sessions~\cite{ref1}.

    \item \textbf{Knowledge Graph-Based Learning Paths:} Constructing a prerequisite knowledge graph over aptitude and coding topics would enable the AI Coach to recommend structured, sequenced learning paths ensuring foundational skills are built before advanced ones are attempted~\cite{ref18}.

    \item \textbf{Multimodal Fusion for Interview Scoring:} Combining audio-derived speech clarity, video-derived presence, and text-derived content quality signals through attention-based multimodal fusion~\cite{ref17} would provide a holistic interview performance metric no single modality captures alone.

    \item \textbf{Personalised Course Recommendations:} Integrating skill gap data from the Resume Analyzer with collaborative filtering over a curated course catalogue would allow the platform to recommend specific online learning resources tailored to each student's identified weaknesses~\cite{ref14}.

    \item \textbf{Competitive Leaderboards and Peer Practice:} Adding an optional cohort leaderboard along with a peer mock interview matching system would introduce social motivation and peer feedback dimensions currently absent.

    \item \textbf{Mobile Application:} A React Native mobile client would extend AceIt to the mobile-first student demographic, enabling on-the-go aptitude practice and GD preparation without requiring a desktop environment.

    \item \textbf{Company-Specific Preparation Packs:} Curating question banks and interview question sets specific to major campus recruiters (TCS, Infosys, Amazon, Google) would allow students to shift from general preparation to targeted, company-specific practice as placement season approaches.

    \item \textbf{Proctored Mock Placement Drives:} Building a full-simulation placement drive mode covering aptitude, coding, GD, and interview rounds in a single proctored session would give students experience managing time and stamina across a complete placement process before facing the real event.
\end{enumerate}

% ============================================================
% REFERENCES
% ============================================================
\begin{thebibliography}{20}

\bibitem{ref1}
S.~Aleven, E.~McLaughlin, R.~Baker, and K.~Koedinger,
``Instruction based on adaptive learning technologies,''
in \textit{Handbook of Research on Learning and Instruction}, 2nd ed., R.~Mayer and P.~Alexander, Eds.
New York, NY, USA: Routledge, 2016, pp.~522--560.

\bibitem{ref2}
A.~Bharadwaj and S.~Saraswat,
``AI-powered resume screening system using NLP and machine learning,''
in \textit{Proc.\ International Conference on Machine Learning and Data Engineering (iCMLDE)}, Sydney, Australia, 2022, pp.~115--122.

\bibitem{ref3}
C.~Ihantola, T.~Ahoniemi, V.~Karavirta, and O.~Seppala,
``Review of recent systems for automatic assessment of programming assignments,''
in \textit{Proc.\ 10th Koli Calling International Conference on Computing Education Research}, Koli, Finland, 2010, pp.~86--93.

\bibitem{ref4}
T.~Shaikh and S.~Gupta,
``Conversational AI for technical interview coaching using large language models,''
\textit{Journal of Educational Technology \& Society}, vol.~27, no.~1, pp.~44--58, 2024.

\bibitem{ref5}
Z.~Li, H.~Zhao, and M.~Chen,
``Real-time automated speech analysis for structured interview feedback,''
\textit{IEEE Transactions on Learning Technologies}, vol.~16, no.~4, pp.~521--532, 2023.

\bibitem{ref6}
P.~Ramesh and K.~Sankar,
``Automated group discussion scoring using multi-dimensional NLP rubrics,''
in \textit{Proc.\ International Conference on Artificial Intelligence in Education (AIED)}, 2023, pp.~302--313.

\bibitem{ref7}
V.~Lugaresi, J.~Tang, H.~Nash, C.~McClanahan, E.~Uboweja, M.~Hays, F.~Zhang, C.~Chang, M.~Yong, J.~Lee, W.~Chang, W.~Hua, M.~Georg, and M.~Grundmann,
``MediaPipe: A framework for building perception pipelines,''
\textit{arXiv preprint arXiv:1906.08172}, 2019.

\bibitem{ref8}
D.~Weiss, N.~Kingsbury, and R.~Betz,
``Adaptive computerized testing: A review,''
\textit{Psychological Bulletin}, vol.~96, no.~2, pp.~339--355, 1984.

\bibitem{ref9}
E.~Alese, F.~Adebiyi, and B.~Omole,
``A survey of cloud-based sandboxed code execution environments for automated assessment,''
\textit{International Journal of Computer Science \& Information Technology}, vol.~15, no.~3, pp.~45--62, 2023.

\bibitem{ref10}
R.~Sharma, A.~Gupta, and P.~Mehta,
``Intelligent resume evaluation tool based on machine learning and natural language processing,''
\textit{International Journal of Advanced Computer Science and Applications (IJACSA)}, vol.~14, no.~2, pp.~201--210, 2023.

\bibitem{ref11}
J.~Wollny, J.~Schneider, D.~Di~Mitri, J.~Weidlich, M.~Rittberger, and H.~Drachsler,
``Are we there yet? -- A systematic literature review on chatbots in education,''
\textit{Frontiers in Artificial Intelligence}, vol.~4, art.\ no.\ 654924, 2021.

\bibitem{ref12}
S.~Geyik, S.~Ambler, and K.~Kenthapadi,
``Fairness-aware ranking in search \& recommendation systems with application to LinkedIn talent search,''
in \textit{Proc.\ ACM SIGKDD International Conference on Knowledge Discovery \& Data Mining}, Anchorage, AK, USA, 2019, pp.~2221--2231.

\bibitem{ref13}
M.~Kazemitabaar, J.~Chow, C.~Ma, B.~Ericson, D.~Weintrop, and T.~Grossman,
``Studying the effect of AI code generators on supporting novice learners in introductory programming,''
in \textit{Proc.\ CHI Conference on Human Factors in Computing Systems}, Hamburg, Germany, 2023, pp.~1--23.

\bibitem{ref14}
S.~Kumar and R.~Verma,
``Intelligent skill gap detection and personalised career guidance system using NLP and collaborative filtering,''
in \textit{Proc.\ International Conference on Advances in Computing, Communication \& Control (ICAC3)}, Mumbai, India, 2024, pp.~1--8.

\bibitem{ref15}
A.~Mehrotra, V.~Sahu, and R.~Pandey,
``ResumAI: Revolutionising automated resume analysis and recommendation with multi-model intelligence,''
\textit{Journal of Artificial Intelligence Research}, vol.~78, pp.~1045--1078, 2023.

\bibitem{ref16}
D.~Ramachandran and M.~Pillan,
``Automated feedback generation for open-ended academic writing using rubric-grounded large language models,''
\textit{Computers \& Education: Artificial Intelligence}, vol.~5, art.\ no.\ 100153, 2023.

\bibitem{ref17}
A.~Suman, T.~Bhatt, and K.~Gupta,
``Multimodal assessment in online education: Fusing audio, video, and text signals for holistic evaluation,''
\textit{IEEE Access}, vol.~11, pp.~72410--72425, 2023.

\bibitem{ref18}
J.~Doroudi, E.~Brunskill, and V.~Aleven,
``Towards a unified theory of state abstraction for MDPs in student modelling,''
in \textit{Proc.\ International Conference on Educational Data Mining (EDM)}, 2019, pp.~324--329.

\bibitem{ref19}
S.~Kumar, R.~Verma, and A.~Nair,
``Real-time AI-based skill gap analysis and adaptive career guidance using Gemini Flash and weighted cosine similarity,''
in \textit{Proc.\ International Conference on Advances in Computing, Communication and Control (ICAC3)}, Mumbai, India, 2024, pp.~1--8.

\bibitem{ref20}
G.~Torres, F.~Ruiz, and P.~Sanchez,
``Using graph algorithms for skills gap analysis: A Neo4j-based approach integrating PageRank, Louvain communities and O*NET ontology,''
\textit{Expert Systems with Applications}, vol.~238, art.\ no.\ 122156, 2024.

\end{thebibliography}

\end{document}
